\documentclass[letterpaper,11pt]{article}
\usepackage[brazil]{babel}
\usepackage[T1]{fontenc}
\usepackage{ae}
\usepackage[utf8]{inputenc}
\usepackage[dvipsnames]{color}
\usepackage{graphicx}
\usepackage{epsfig}
\usepackage{multirow}
\usepackage{amssymb, amsmath, amsfonts}
\bibliographystyle{plain}
\usepackage{url}
\topmargin		0 cm
\hoffset		0 cm
\voffset		0 cm
\evensidemargin		0 cm
\oddsidemargin		0 cm
\setlength{\textwidth}{16 cm}
\setlength{\textheight}{21 cm}
\usepackage{../../tutoriais}

\title{Apostila de Instalação do CUTEst}
\author{
Abel Soares Siqueira\footnote{abel.s.siqueira@gmail.com} \and
Douglas Gonçalves\footnote{douglas@ime.unicamp.br - desatualizado} \and
Leandro F. Prudente\footnote{lfprudente@ime.unicamp.br - desatualizado} \and
Raniere Gaia Costa da Silva\footnote{ra092767@ime.unicamp.br}}

\numberwithin{equation}{section}

% \newcommand{\visiblelbl}[1]{\label{#1}{\color{red}{#1}}}
\newcommand{\visiblelbl}[1]{\label{#1}}

\begin{document}
\maketitle
\tableofcontents
\vspace{1 cm}

A biblioteca CUTEst (\cite{bib:cuter1,bib:cuter2,bib:cutest}) é de extrema
importância para a realização de testes computacionais na área de otimização
computacional.  Tendo em vista que é necessário fazer comparações entre os
algoritmos, Muitos dos mais famosos atualmente têm uma interface que permite a
execução dos testes do CUTEst. Vamos mostrar nesse tutorial como instalar o
CUTEst e criar uma interface em Fortran e C para o mesmo. Algumas partes deste
tutorial supõem que o leitor saiba o suficiente de linux.

Para uma interface em Windows ou Mac, leia veja a Seção \ref{sec:contribute}.

%!TEX root = cutest.tex
\section{Introdução}
% \subsection{Como ler este tutorial}
% \subsection{Autores e Agradecimentos}
\subsection{Obtendo ajuda}
Quando você precisar de ajuda com assuntos relacionados a este tutorial, faça o seguinte
\begin{enumerate}
    \item Procure em \url{https://github.com/abelsiqueira/tutoriais/issues} pelo
        seu problema. Provável que alguém tenha reportado um erro similar.
    \item Dê uma lida na documentação oficial do CUTEst pois pode ser que este
        tutorial esteja desatualizado.
    \item Entre em contato com um dos autores.
    \item Procure no irc, servidor \verb+freenode+ canal \verb+##cutest+.
\end{enumerate}
Tenha em mente que estamos fazendo esse serviço de bondade. 
Não temos nenhuma relação direta com o projeto ou seus autores. 
Nosso conhecimento sobre o assunto é limitado.
\section{Instalação}
O CUTEr é um ambiente de testes para métodos de otimização. 
Ele foi feito em Fortran, porém tem suporte a C/C++.
Para instalar o CUTEr você irá precisar dos programas:
\begin{itemize}
 \item gawk,
 \item gcc,
 \item gfortran (ou outro compilador de fortran),
 \item svn (subversion).
\end{itemize}
Com os pacotes instalados, iremos criar uma pasta para o CUTEr. Você pode criar uma
pasta separada no sistema ou instalar tudo no seu diretório principal. No meu caso,
irei criar uma pasta chamada \verb+libraries+ na minha pasta pessoal, e criar uma
pasta para o CUTEr dentro dessa pasta. Ajuste os comandos de acordo com sua escolha.
\begin{terminal}
$ mkdir -p $HOME/libraries/CUTEr
$ cd $HOME/libraries/CUTEr
\end{terminal}
Agora faça o download do SifDec2 pelo comando a seguir
\begin{terminal}
$ svn co http://tracsvn.mathappl.polymtl.ca/SVN/cuter/sifdec/branches/SifDec2 ./sifdec2
\end{terminal}
E você precisa baixar o CUTEr2 agora. Se quiser a versão 32 bits:
\begin{terminal}
$ svn co http://tracsvn.mathappl.polymtl.ca/SVN/cuter/cuter/branches/CUTEr2 ./cuter2
\end{terminal}
E a versão 64 bits:
\begin{terminal}
$ svn co http://tracsvn.mathappl.polymtl.ca/SVN/cuter/cuter/branches/cuter64 ./cuter2
\end{terminal}
Faça o download da coleção de problemas SIF na página
\url{http://cuter.rl.ac.uk/cuter-www/Problems/mastsif.shtml} ou use os comandos
\begin{terminal}
$ wget ftp://ftp.numerical.rl.ac.uk/pub/cuter/mastsif_small.tar.gz 
$ wget ftp://ftp.numerical.rl.ac.uk/pub/cuter/mastsif_large.tar.gz
\end{terminal}
Descompacte os problemas pequenos com o comando
\begin{terminal}
$ tar -zxvf mastsif_small.tar.gz 
\end{terminal}
Se quiser, descompacte os problemas grandes também. Note que estes problemas são
muito mais pesados que os pequenos, então aconselho baixá-los apenas se for necessário.
\begin{terminal}
$ tar -zxvf mastsif_large.tar.gz
\end{terminal}
Iremos começar a instalar o CUTEr. Para isso, vamos instalar o decodificador
dos problemas SIF (que faz parte da biblioteca CUTEr).
Entre na pasta \verb+sifdec2+ criada em \verb+$HOME/libraries/CUTEr+.
\begin{terminal}
$ cd $HOME/libraries/CUTEr/sifdec2
\end{terminal}
Agora digite o comando
\begin{terminal}
$ ./install_sifdec
\end{terminal}
A instalação do sifdec irá pedir que você escolha
\begin{itemize}
 \item Plataforma (no nosso caso será 5 - PC)
 \item Sistema Operacional (no nosso caso será 2 - linux)
 \item Compilador Fortran (no nosso caso será 7 - GNU gfortran)
 \item Precisão (no nosso caso será D - double)
 \item Tamanho (no nosso caso será L - large, porém pode ser necessário mudar para C - custom,
dependendo do software e dos testes que você utilizará. Nesse caso edite o arquivo na pasta
\verb+build/arch/size.custom+ antes de continuar.)
\end{itemize}
O instalador vai lançar a mensagem
\begin{terminal}
 By default, SifDec with your selections will be installed in
  /home/abel/libraries/CUTEr/sifdec2/SifDec.large.pc.lnx.gfo
 Is this OK (Y/n)?
\end{terminal}
Anote o nome do diretório que será criado e aperte enter. No nosso caso, será \\
 \verb+SifDec.large.pc.lnx.gfo+. Para a próxima mensagem, escreva `n' e aperte enter.

Agora vamos acrescentar as variáveis de sistema para que o instalador e os softwares
que formos utilizar saibam onde está o CUTEr. Edite o arquivo \verb+$HOME/.bashrc+ e
adicione as seguintes linhas
\begin{terminal}
ROOTCUTER="$HOME/libraries/CUTEr"
export CUTER="$ROOTCUTER/cuter2"
export MYCUTER="$CUTER/CUTEr.large.pc.lnx.gfo" 
export SIFDEC="$ROOTCUTER/sifdec2"
export MYSIFDEC="$SIFDEC/SifDec.large.pc.lnx.gfo" 
export MASTSIF="$ROOTCUTER/mastsif" 
export MANPATH="$CUTER/common/man:$SIFDEC/common/man:$MANPATH" 
export PATH="$MYCUTER/bin:$MYSIFDEC/bin:$PATH"
\end{terminal}
Se for necessário, mude a variável \verb+ROOTCUTER+, e o diretório que eu mandei anotar. 
Note que já criamos um diretório para o CUTEr, supondo que ele siga as mesmas opções
que o SifDec. 
Para que essas mudanças façam efeito é necessário usar o comando
\begin{terminal}
$ source $HOME/.bashrc
\end{terminal}
ou reiniciar o terminal.
Agora entre na pasta que criamos anteriormente. Para isso basta fazer
\begin{terminal}
$ cd $MYSIFDEC
\end{terminal}
Se deu algum erro, você pode ter errado o nome na definição da variável. Verifique que a
pasta realmente existe e refaça os passos se necessário.

Agora rode o comando
\begin{terminal}
$ ./install_mysifdec
\end{terminal}
Ao final desse comando, a frase
\begin{terminal}
install_mysifdec : Do you want to 'make all' in
 /home/abel/libraries/CUTEr/sifdec2/SifDec.large.pc.lnx.gfo now (Y/n)?
\end{terminal}
aparecerá. Aperte enter. Agora vamos instalar o CUTEr. Para isso, faça
\begin{terminal}
$ cd $CUTER
$ ./install_cuter
\end{terminal}
A instalação do CUTEr irá pedir que você escolha
\begin{itemize}
 \item Plataforma (no nosso caso será 5 - PC)
 \item Sistema Operacional (no nosso caso será 2 - linux)
 \item Compilador Fortran (no nosso caso será 7 - GNU gfortran)
 \item Compilador C (no nosso caso será 4 - g++)
 \item Precisão (no nosso caso será D - double)
 \item Tamanho (no nosso caso será L - large, com a mesma observação sobre o tamanho.)
\end{itemize}
Quando necessário, aperte enter para confirmar a continuação do programa. Se tudo deu
certo, o CUTEr deve estar instalado. Caso tenha acontecido algum erro, reveja as variáveis
e volte os passos. Pode ser necessário apagar tudo que você baixou e baixar novamente (em 
último caso). Para verificar que realmente está tudo funcionando, faça
\begin{terminal}
$ mkdir -p $HOME/libraries/CUTEr/Testing
$ cd $HOME/libraries/CUTEr/Testing
$ runcuter -p gen -D ROSENBR
\end{terminal}
Se nenhum erro aparecer, e aparecer várias informações do problema, então o
CUTEr foi instalado corretamente.

\section{Interface}
A biblioteca CUTEr dá ao usuário um conjunto de funções para receber informações do problema.
As funções são separadas para o caso irrestrito (precedidas pela letra \emph{U}) ou
restrito, (precedidas pela letra \emph{C}). Se o usuário necessitar do valor da funcão
 objetivo num ponto dado, ele pode chamar \verb+UFN+ se o problema for irrestrito, ou
\verb+CFN+ caso contrário. Note que a sintaxe das funções não é a mesma, pois o CUTEr tenta
ser o mais prático possível. A sintaxe dessas funcões é
\begin{itemize}
 \item \verb+UFN (N, X, F)+, onde \verb+N+ é um inteiro indicando o número de 
variáveis do problema, \verb+X+ é um vetor de
reais, e \verb+F+ é real que sai com o valor da função objetivo.
 \item \verb+CFN (N, M, X, F, LC, C)+, onde \verb+N+,\verb+X+ e \verb+F+ indicam as
mesmas coisas, \verb+M+ é um inteiro indicando o número de restrições, \verb+C+ é um
vetor de reais, com os valores das restrições e \verb+LC+ é a dimensão real de \verb+C+,
que deve ser maior ou igual a \verb+M+.
\end{itemize}
Existem muitas outras funções para interação com o problema. Elas podem ser vistas na
documentação do CUTEr no arquivo \verb+general.pdf+ dentro da pasta
\verb+$CUTER/common/doc+ nas páginas 31 e 32.

Para rodar nossa biblioteca com o CUTEr, é necessário criar alguns arquivos extras.
Um deles é o código que relaciona o nosso programa com essas funções do CUTEr. Além disso,
também é necessário criar um arquivo que indica ao CUTEr quais as bibliotecas que nossa
biblioteca precisa. A maneira tradicional de se trabalhar com o CUTEr (o novo CUTEr, pelo
menos) é compilar nossa biblioteca inteira para um arquivo .a, compilar esse arquivo de
ligação do CUTEr com nossa biblioteca, e passar tudo isso pro CUTEr. Quando necessário,
rodamos o nosso pacote pelo comando do CUTEr \verb+runcuter+. Por exemplo, para rodar o
pacote \verb+pack_exemplo+, usamos o comandos
\begin{terminal}
$ runcuter -p pack_exemplo -D problema
\end{terminal}
onde \verb+problema+ é um dos problemas do CUTEr sem a extensão \verb+.SIF+.

\subsection{Instalação da interface}

\section{Exemplos}
Como exemplos do CUTEr, vamos criar bibliotecas em diversas linguagens para resolver o 
problema de minimização irrestrita. 
Vamos implementar o método de máxima descida com busca linear utilizando o
critério de Armijo.

Considere o problema
\begin{equation*}
 \min f(x), \qquad x \in \Rn{n},
\end{equation*}
onde $f:\RnemR{n}$ é contínua e derivável. Vamos procurar um ponto estacionário para esse
problema, isto é, um ponto $x^* \in \Rn{n}$ tal que $$\nabla f(x^*) = 0.$$
Obviamente, como vamos implementar este método, vamos parar quando encontrarmos um
iterando $x^k$ tal que $\norma{\nabla f(x^k)} \leq \varepsilon,$ onde $\varepsilon > 0$ é
dado. O método está descrito a seguir.
\begin{algorithm}[H]
\caption{Método de Máxima Descida}
 \begin{algorithmic}[1]
 \State Dados $x^0 \in \Rn{n}$, $\varepsilon > 0$, $\alpha \in (0,1)$, $k = 0$.
 \While{$\norma{\nabla f(x^k)} > \varepsilon$}
  \State $d^k = -\nabla f(x^k)$
  \State $\lambda_k = 1$
  \While {$f(x^k + \lambda_kd^k) > f(x^k) + \alpha \lambda_k \nabla f(x^k)^Td^k$}
   \State $\lambda_k = \lambda_k/2$
  \EndWhile
  \State $x^{k+1} = x^k + \lambda_kd^k$
  \State $k = k + 1$
 \EndWhile
 \State $x^* = x^k$.
 \end{algorithmic}
\end{algorithm}
Vamos implementar este método em algumas linguagens, e às vezes, mais de uma vezes, para
exemplificar a interface CUTEr.

\subsection{Instruções Gerais}

\emph{
Os passos a seguir são feitos normalmente após a criação da biblioteca, porque
primeiro você precisa fazer com que ela funcione em outros problemas. Na dúvida,
faça um dos exemplos a seguir e depois volte para ler esta seção.}

Para que seu programa funcione com o CUTEr, você irá precisar de um arquivo de 
protótipo que pode ser copiado de \verb+$CUTER/build/prototypes/gen90.sh.pro+.
Nesse arquivo, você irá mudar algumas variáveis para o compilador do CUTEr
encontrar seus arquivos. As linhas sem comentários estão mostradas a seguir:
\begin{code}{bash}
export PACK=gen90
export PACKAGE=gen
export PACK_PRECISION="single double"
export PACKOBJS="cuterinter.o gen90.o"
export PACKLIBS=""
export SPECS="GEN.SPC"
\end{code}
Vamos supor que temos um pacote chamado \verb+cutertest+, e que esse pacote precisa
de um objeto \verb+auxiliar.o+, e de uma biblioteca do sistema \verb+libalgo.a+.
Além disso, você deve ter criado (veja nos exemplos) um arquivo específico para
fazer a interface entre o CUTEr e o pacote \verb+cutertest+. Vamos supor que
a linguagem em questão seja \verb+fortran+, então você deve ter criado um
arquivo \verb+cutertestmain.f+. A variável \verb+PACK+ define um pequeno
nome para o seu pacote, e \verb+PACKAGE+ é o nome de seu pacote.
\verb+PACK_PRECISION+ é a precisão do seu pacote (muitas vezes double).
\verb+PACKOBJS+ são os objetos que você deve colocar na pasta
\verb+$MYCUTER/(precision)/bin+ onde (precision) é a precisão que você instalou
o CUTEr.
\verb+PACKLIBS+ são as bibliotecas que o pacote precisa, e \verb+SPECS+ se refere
a algum arquivo de especificações que o seu pacote pode depender. Esse arquivo pode
ser local (isto é, na pasta onde você irá chamar o CUTEr) ou na pasta
\verb+$CUTER/common/src/pkg/$PACKAGE/+.

Para o nosso exemplo deveríamos ter no arquivo chamado \verb+cutertest.sh.pro+
as seguintes definições:
\begin{code}{bash}
export PACK=cutertest
export PACKAGE=cutertest
export PACK_PRECISION="double"
export PACKOBJS="auxiliar.o"
export PACKLIBS="-lalgo"
export SPECS=""
\end{code}
Note que devemos copiar o arquivo \verb+auxiliar.o+ para \verb+$MYCUTER/double/bin+.
Agora, precisamos passar esse arquivo por um formatador de texto utilizando um
script que o CUTEr fornece e devemos habilitá-lo para execução. Para isso, os comandos são
\begin{terminal}
$ sed -f $MYCUTER/double/config/script.sed cutertest.sh.pro > $MYCUTER/bin/cutertest
$ chmod a+x $MYCUTER/bin/cutertest
\end{terminal}
Agora passamos o arquivo de interface e o arquivo compilado para o CUTEr.
\begin{terminal}
$ cp -f cutertestmain.f $MYCUTER/double/bin/cutertestma.f
$ cp -f cutertestmain.o $MYCUTER/double/bin/cutertestma.o
\end{terminal}
Tendo feito isso, agora seu arquivo está pronto para ser executado. Use o comando
\begin{terminal}
$ runcuter -p cutertest -D ROSENBR
\end{terminal}

%!TEX root=cutest.tex
\subsection{Exemplos em Fortran}
Fizemos uma implementação do método de máxima descida. Temos dois arquivos na
implementação:
\begin{itemize}
 \item \verb+gradient.f+: Este arquivo contém a definição do método.
 \item \verb+gradientmain.f+: Este arquivo contém a rotina principal do fortran.
\end{itemize} Além desses arquivos também é necessário um arquivo com as
definições das subrotinas
\begin{itemize}
 \item \verb+inip(n,x)+: Retorna n, a dimensão do problema, e x, o ponto
 inicial.
 \item \verb+evalf(n,x,f)+: Recebe a dimensão do problema n, e o ponto x, e
retorna o valor da função objetivo em f.
 \item \verb+evalg(n,x,g)+: Recebe a dimensão do problema n, e o ponto x, e
retorna o valor do gradiente da função objetivo em g.
 \item \verb+endp(n,x)+: Imprime informações sobre a solução.
\end{itemize}
Para criar a interface em fortran, é necessário apenas criar um arquivo com as
subrotinas acima. O arquivo com a interface (sem os comentários está abaixo:)
\lstinputlisting[style=codestyle, language=fortran]
{../Examples/FortranExample1/fortran_examplemain.f}

%!TEX root=cutest.tex
\subsection{Exemplos em C}
Fizemos três implementações do método de máxima descida. A primeira é uma implementação
que não leva me conta o CUTEr, e depois adapta o CUTEr para o problema. A segunda já leva
em conta o formato das funções do CUTEr e faz pouca adaptação posteriormente. A terceira
usa exatamente as funções do CUTEr, não necessitando de adaptação.

Cada implementação do método de máxima descida consiste de dois arquivos:
\verb+steepest_descent.h+ e \verb+steepest_descent.c+. No .h, definimos que funções
iremos chamar, e uma estrutura com as informações da execução. As funções para a 
primeira implementação são
\begin{itemize}
 \item \verb+double Norm (double * x, unsigned int n);+ \\
Esta função calcula a norma 2 de um vetor \verb+x+ com tamanho \verb+n+.
 \item \verb+double NormSqr (double * x, unsigned int n);+ \\
Esta função calcula o quadrado da norma 2 de um vetor \verb+x+ com tamanho \verb+n+. É mais
rápido que a função \verb+Norm+ pois não envolve raiz quadrada.
 \item \verb+SteepestDescent (double * x, unsigned int n, Status * status);+ \\
Esta é a função que encontra o ponto estacionário. \verb+x+ entra como ponto inicial e 
sai como a solução. \verb+n+ é a dimensão do problema e \verb+status+ é um ponteiro para
a estrutura de informações.
 \item \verb+SD_Print (double * x, unsigned int n, Status * status);+ \\
Esta função imprime o vetor \verb+x+ e as informações da execução do problema.
\end{itemize}
A estrutura do problema
\begin{code}{C}
typedef struct _Status {
  unsigned int iter;
  double f, ng;
  unsigned int n_objfun, n_gradfun;
} Status;
\end{code}
\verb+iter+ é o número de iterações que o algoritmo executou, \verb+f+ é o valor da
função objetivo na solução, \verb+ng+ é a norma do gradiente da função objetivo,
\verb+n_objfun+ é o número de cálculos da função objetivo e \verb+n_gradfun+ é o número
de cálculos do gradiente. Mostraremos as diferenças das outras implementações posteriormente.

Nosso arquivo .c contém as definições das funções acima, e contém uma declaração de função
usada para acessar a função objetivo e o gradiente. Na primeira implementação, essa declaração
é
\begin{code}{C}
double objfun  (double * x, unsigned int n);
void   gradfun (double * x, unsigned int n, double * g);
\end{code}
As funções \verb+Norm+, \verb+NormSqr+ e \verb+SD_print+ são idênticas para todas as
implementações e serão deixadas de fora. A implementação do método em si encontra-se abaixo.
\scriptsize
\lstinputlisting[style=codestyle, language=C]
{../Examples/CExample1/steepest_descent.c}
\normalsize
Um código de exemplo para esse teste é
\scriptsize
\lstinputlisting[style=codestyle, language=C]
{../Examples/CExample1/test1.c}
\normalsize
Este código implementa o problema de minimizar $f(x) = \meio\norma{x}^2$ em duas dimensões.
Note como temos que declarar as funções \verb+objfun+ e \verb+gradfun+. Sem elas teríamos
erros na compilação. Veja os arquivos \verb+test2.c+ e \verb+test3.c+
para outros exemplos.

A interface para o CUTEr é o arquivo \verb+c_example1main.c+:
\scriptsize
\lstinputlisting[style=codestyle, language=C]
{../Examples/CExample1/c_example1main.c}
\normalsize
Note que também é necessário definir as funções \verb+objfun+ e \verb+gradfun+.
Essas funções, por sua vez, chamam as funções correspondentes em CUTEr para esse
serviço. A função \verb+UFN+ calcula o valor da função objetivo e a função
\verb+UGR+ calcula o valor do gradiente. Note que quando compilamos uma função
em Fortran, ele recebe um nome em letras minúsculas e com um \_ (underline) na
frente (no caso do gfortran. Outros compiladores podem divergir). O arquivo
\verb+cuter.h+ define macros para todas as funções do CUTEr serem chamados com
letras maiúsculas. Note ainda que a função \verb+UFN+ recebe um ponteiro para
\verb+int+ e dois ponteiros para \verb+double+, sendo o primeiro para o vetor
$x$ e o segundo para o valor da função objetivo. As funções do CUTEr para C
recebem ponteiros em todos os valores. Os valores que são vetores não precisam
de um ponteiro adicional. Note também que como utilizamos \verb+unsigned int+
para os tamanhos, tivemos que converter os valores para \verb+int+.

Veja agora o Exemplo 2. Nesse exemplo consideramos que as funções devem estar no mesmo
formato que a função do CUTEr.
\begin{code}{C}
void ufn (int * n, double * x, double * f);
void uofg (int * n, double * x, double * f, double * g, long int * grad);
\end{code}
O nome das funções foram escolhidas para seguir exatamente o formato do CUTEr, mas aqui
elas poderiam ser qualquer coisa. Note que a função \verb+uofg+ foi utilizada no lugar
da função \verb+ugr+. Essa função já calcula a função objetivo e o gradiente, sendo mais
rápida que chamadas individuais. Com essa mudança, esse programa é levemente mais rápido
que o outro. O resto do arquivo foi mudado de acordo a seguir essas mudanças.

A interface também teve uma mudança na definição das funções:
\begin{code}{C}
void ufn (int * n, double * x, double * f) {
  UFN(n, x, f);
}

void uofg (int * n, double * x, double * f, double * g, long int * grad) {
  UOFG(n, x, f, g, grad);
}
\end{code}
Essa interface fica muito mais natural de ser utilizada num problema com CUTEr. Não precisamos
converter nenhuma variável, simplesmente fazer a chamada da função com os parâmetros já dados.
Note, no entanto, que estamos utilizando \verb+long int+ para as variáveis lógicas do CUTEr.
Essa é uma definição do arquivo \verb+cuter.h+. Se mudarmos essa definição, então devemos
criar uma variável \verb+logical GRAD = *grad+ e chamar a função com
\verb+UOFG(n, x, f, g, &GRAD)+.

A última interface já declara as funções que o CUTEr irá definir. Então no arquivo .c temos
\begin{code}{C}
void ufn_ (int * n, double * x, double * f);
void uofg_ (int * n, double * x, double * f, double * g, long int * grad);
\end{code}
e no arquivo da interface não temos nenhuma definição de função. Lembre-se que o Fortran
compilado com gfortran cria as funções com minúsculas e \_ na frente. Se mudarmos o
compilador pode não funcionar. 
Além disso, a definição dessas funções é, na verdade
\begin{code}{Fortran}
void UFN( integer *n, doublereal *x, doublereal *f );
void UOFG( integer *n, doublereal *x, doublereal *f, doublereal *g,
      logical *grad);
\end{code}
e esses tipos são definidos no arquivo \verb+cuter.h+, podendo ser alterados pelo usuário.
Se isso acontecer, é necessário mudar todo o programa.

Uma alternativa é utilizar \verb+typedef+s para definir os tipos próprios, e deixar o
acesso desses tipos para o usuário (assim como o arquivo \verb+cuter.h+). Dessa maneira,
se o usuário tiver necessidade de mudar o arquivo \verb+cuter.h+, ele também pode (deve)
mudar o arquivo com esses \verb+typedef+s.

Para compilar a interface CUTEr dos testes utilize
\begin{terminal}
$ make cuter
\end{terminal}
Para rodar os testes utilize o comando
\begin{terminal}
$ runcuter -p c_example# -D BARD
\end{terminal}
onde \# é o número do exemplo e BARD é um dos problemas em que esse exemplo converge. Note
que é preferível criar uma pasta separada para rodar os testes, já que eles geram lixo na
pasta.

% \subsection{Exemplos em C++}

% \subsection{Exemplos em Python}

\subsection{Exemplos em MATLAB}
Com o {\tt CUTEr} devidamente instalado, vejamos como ocorre a interface com o {\tt MATLAB}.

Primeiro, crie uma pasta {\tt Test}, onde ser\~ao gerados os arquivos decodificados do problema
 e o arquivo {\tt .mex} da interface para {\tt MATLAB}.

No terminal, dentro da pasta {\tt Test}, execute o comando:
\begin{terminal}
$ runcuter -p mx -D  ROSENBR
\end{terminal}
Al\'em dos arquivos padr\~ao, gerados pelo decoder para o problema {\tt ROSENBR}, haver\'a 
tamb\'em o arquivo: {\tt mcuter.mex}. A extens\~ao {\tt .mex} poder\'a variar conforme a 
instala\c{c}\~ao do {\tt CUTEr} e o sistema operacional.

Em seguida, v\'a para o {\tt MATLAB}. Adicione as pastas {\tt \$CUTER/common/src/matlab} e {\tt 
\$MYCUTER/bin} ao {\tt path} do {\tt MATLAB}. Isso pode ser feito atrav\'es do menu {\tt File\ 
>\ Set Path...} ou pelo prompt do {\tt MATLAB} usando o comando {\tt addpath}.

Dentro do {\tt MATLAB}, v\'a para a pasta {\tt Test}. Para verificar se est\'a tudo certo, 
execute o comando:
\begin{code}{matlab}
>> prob = cuter_setup()

prob = 

         n: 2
         m: 0
      nnzh: 3
      nnzj: 0
         x: [2x1 double]
        bl: [2x1 double]
        bu: [2x1 double]
         v: [0x1 double]
        cl: [0x1 double]
        cu: [0x1 double]
    equatn: [0x1 logical]
    linear: [0x1 logical]
      name: 'ROSENBR   '
\end{code}
Se estiver tudo correto, a sa\'ida dever\'a ser como acima.

O comando {\tt cuter\_setup}, que inicializa o problema e fornece suas carcacter\'isticas, se 
encontra na pasta {\tt \$CUTER/src/common/matlab}, dentro da qual est\~ao as demais rotinas 
(arquivos {\tt .m}) para acessar fun\c{c}\~ao objetivo, gradiente, restri\c{c}\~oes, etc.

Por exemplo, para saber o valor da fun\c{c}\~ao objetivo em um determinado ponto, usamos:
\begin{code}{matlab}
>> f = cuter_obj([-1 2]')

f =

   104
\end{code}
Abaixo temos um c\'odigo em {\tt MATLAB} que implementa o m\'etodo do gradiente. Para acessar a 
fun\c{c}\~ao objetivo e o gradiente, fizemos uso da fun\c{c}\~ao {\tt cuter\_obj}.
\begin{code}{matlab}
function [x,f,gradnorm,iter,flag] = cute_gradient()
%
%   [x,f,gradnorm,iter,flag] = cute_gradient()
%
%   flags:
%   -2      maximum number of iterations reached
%   -1      too short step size
%    1      gradient norm less than eps0
%

    % inicializando o problema
    prob = cuter_setup();
    
    % dimensao do problema
    n = prob.n;
    
    % ponto inicial
    x0 = prob.x;
    
    %====================================
    % Gradient method with linesearch
    %====================================
    maxit=n*10000;
    gamma=1e-4;
    eps0=1e-5;
    tmin=1e-8;
    done=0;
    flag=0;
    k=0;
    x=x0;
    
    while ~done
        k=k+1;
        
        % maximum number of iterations test
        if k>maxit
           done=1;
           flag=-2;
           continue
        end
        
        x0=x;
        
        [f,g] = cuter_obj(x);
        f0=f;
        
        gradnorm=norm(g,2);
        
        % gradient norm test
        if gradnorm<eps0
            done=1;
            flag=1;
            continue
        end
        
        % search direction
        d = -g;
        gtd = g'*d;
        t = 1;
        
        x = x0 + t*d;
        f = cuter_obj(x);
        
        % Armijo linesearch
        while f > f0 + t*gamma*gtd
            t = 0.5*t;
            if t<tmin
                done=1;
                flag=-1;
            end
            
            x = x0 + t*d;
            f = cuter_obj(x);
        end
        
    end
    
    iter=k;

end
\end{code}
E executando o c\'odigo acima, obtemos
\begin{code}{matlab}
>> [x,f,gradnorm,iter,flag] = cute_gradient()

x =

    1.0000
    1.0000


f =

   6.1319e-11


gradnorm =

   9.9971e-06


iter =

       10917


flag =

     1
\end{code}

\'E poss\'ivel tamb\'em utilizar as rotinas de otimiza\c{c}\~ao do pr\'oprio {\tt MATLAB}. Por 
exemplo, para utilizar o {\tt fminunc}, usamos:
\begin{code}{matlab}
>> prob = cuter_setup();
>> [x,f,flag] = fminunc(@(x) cuter_obj(x),prob.x)
Warning: Gradient must be provided for trust-region method;
  using line-search method instead. 
> In fminunc at 356

Local minimum found.

Optimization completed because the size of the gradient is less than
the default value of the function tolerance.

<stopping criteria details>


x =

    1.0000
    1.0000


f =

   2.8336e-11


flag =

     1
\end{code}
Lembrando que as op\c{c}\~oes do solver s\~ao ajustadas pelo comando {\tt optimset}.
\begin{code}{matlab}
>> prob = cuter_setup();
>> opts = optimset('GradObj','on');
>> [x,f,flag] = fminunc(@(x) cuter_obj(x),prob.x,opts)

Local minimum possible.

fminunc stopped because the final change in function value relative to 
its initial value is less than the default value of the function tolerance.

<stopping criteria details>


x =

    1.0000
    1.0000


f =

   4.0035e-13


flag =

     3
\end{code}


%!TEX root = cutest.tex
\section{Contribua}
\label{sec:contribute}

Contribua

\appendix
\section{Licença}
Esta obra está licenciada sob a licença Creative Commons Atribuição 3.0 Não
Adaptada. Para ver uma cópia desta licença, visite
\url{http://creativecommons.org/licenses/by/3.0/}.
\begin{center}
    \includegraphics{../../figures/cc-by.png}
\end{center}


\begin{thebibliography}{3}
\bibitem{bib:cuter1} http://cuter.rl.ac.uk/cuter-www
\bibitem{bib:cuter2} https://magi-trac-svn.mathappl.polymtl.ca/Trac/cuter

\end{thebibliography}
\end{document}
